\documentclass[11pt, a4paper]{article}

% ---------- Packages ----------
\usepackage[utf8]{inputenc}
\usepackage[T1]{fontenc}
\usepackage[margin=1in]{geometry}
\usepackage{graphicx}
\usepackage{amsmath, amssymb}
\usepackage{booktabs}
\usepackage{caption}
\usepackage{subcaption}
\usepackage{hyperref}
\usepackage[round]{natbib}
\usepackage{siunitx}
\usepackage{float}
\usepackage{fancyhdr}

\hypersetup{
    colorlinks=true,
    linkcolor=black,
    citecolor=blue,
    urlcolor=blue
}

\pagestyle{fancy}
\fancyhf{}
\fancyhead[L]{\nouppercase{\leftmark}}
\fancyfoot[C]{\thepage}

% ---------- Title information ----------
\title{Report Title}
\author{
    First Author\thanks{Corresponding author: \texttt{author@example.com}} \\
    Department, Institution
    \and
    Second Author \\
    Department, Institution
}
\date{\today}

\begin{document}

\maketitle

\begin{abstract}
Summarize the motivation, methods, key results, and conclusions of the
report in 150--250 words. The abstract should be self-contained and
understandable without reading the rest of the document.
\end{abstract}

\noindent\textbf{Keywords:} keyword one, keyword two, keyword three

\tableofcontents
\newpage

% ---------- Introduction ----------
\section{Introduction}
\label{sec:introduction}

State the problem, its significance, and the goal of this report. Summarize
prior work and clearly identify the gap this report addresses. End the
section with an outline of the remaining sections.

% ---------- Related Work ----------
\section{Related Work}
\label{sec:related-work}

Discuss relevant prior studies and how the present work relates to or
differs from them \citep{example2024}.

% ---------- Methods ----------
\section{Methods}
\label{sec:methods}

Describe the materials, experimental setup, data sources, and analytical or
computational methods used. Provide enough detail for reproducibility.

\subsection{Experimental Setup}
Describe the apparatus, dataset, or system configuration.

\subsection{Data Analysis}
Describe statistical tests, models, or algorithms applied to the data.
For example, a variable may be defined as:
\begin{equation}
    \label{eq:example}
    y = f(x) = \beta_0 + \beta_1 x + \varepsilon
\end{equation}
where $\beta_0$ and $\beta_1$ are coefficients and $\varepsilon$ is the
error term.

% ---------- Results ----------
\section{Results}
\label{sec:results}

Present the findings using figures and tables. Report quantitative results
with appropriate units and uncertainty.

\begin{figure}[H]
    \centering
    % \includegraphics[width=0.8\textwidth]{figures/example-figure.pdf}
    \caption{Description of the figure, including axes, units, and what
    the reader should take away.}
    \label{fig:example}
\end{figure}

\begin{table}[H]
    \centering
    \caption{Summary of results.}
    \label{tab:example}
    \begin{tabular}{lcc}
        \toprule
        Condition & Metric A & Metric B \\
        \midrule
        Baseline  & 0.00     & 0.00     \\
        Treatment & 0.00     & 0.00     \\
        \bottomrule
    \end{tabular}
\end{table}

Reference results explicitly, e.g., as shown in Figure~\ref{fig:example}
and Table~\ref{tab:example}, and in Equation~\ref{eq:example}.

% ---------- Discussion ----------
\section{Discussion}
\label{sec:discussion}

Interpret the results in the context of the stated goal and related work.
Discuss limitations, potential sources of error, and alternative
explanations.

% ---------- Conclusion ----------
\section{Conclusion}
\label{sec:conclusion}

Summarize the main findings and their implications. Suggest directions for
future work.

% ---------- Acknowledgments ----------
\section*{Acknowledgments}
Acknowledge funding sources, contributors, and any assistance received.

% ---------- Bibliography ----------
\bibliographystyle{plainnat}
\bibliography{references}

% ---------- Appendix ----------
\appendix
\section{Supplementary Material}
\label{sec:appendix}

Include supplementary derivations, additional figures/tables, or raw data
summaries here.

\end{document}
